% Paquetes básicos
\usepackage[utf8]{inputenc}
\usepackage[spanish]{babel}
\usepackage[T1]{fontenc}
\usepackage{lmodern}

% Geometría y márgenes
\usepackage[a4paper, margin=2.5cm]{geometry}

% Matemáticas
\usepackage{amsmath, amssymb, amsthm}

% Gráficos y figuras
\usepackage{graphicx}
\usepackage{float}
\usepackage{tikz}
\usepackage{pgfplots}
\pgfplotsset{compat=1.18}

% Código fuente
\usepackage{listings}
\usepackage{xcolor}

% Enlaces e hipervínculos
\usepackage[colorlinks=true, linkcolor=blue, urlcolor=blue, citecolor=blue]{hyperref}

% Bibliografía
\usepackage{cite}

% Configuración de listings para C++
\lstset{
    language=C++,
    basicstyle=\ttfamily\footnotesize,
    keywordstyle=\color{blue}\bfseries,
    commentstyle=\color{green!60!black},
    stringstyle=\color{red},
    numbers=left,
    numberstyle=\tiny\color{gray},
    stepnumber=1,
    numbersep=5pt,
    backgroundcolor=\color{gray!10},
    frame=single,
    tabsize=4,
    showspaces=false,
    showstringspaces=false,
    breaklines=true
}

% Definir teoremas
\newtheorem{theorem}{Teorema}[chapter]
\newtheorem{definition}{Definición}[chapter]
\newtheorem{example}{Ejemplo}[chapter]
